% from ega2-indexes-and-contents-fr.tex
\oldpage[II]{209}
\markboth{A. GROTHENDIECK --- Chap. II}{A. GROTHENDIECK --- Chap. II}
\thispagestyle{plain}
\phantomsection
\addcontentsline{toc}{section}{Index terminologique}
\section*{Index terminologique}

\begingroup
\small
\setlength{\parindent}{0pt}
\setlength{\parskip}{2pt}
\newcommand{\egaTerme}[2]{\textbf{#1} : #2.\par}

\egaTerme{Affine (morphisme)}{\hyperref[II.1.6.1-fr]{1.6.1}}
\egaTerme{Affine (préschéma) au-dessus d'un préschéma}{\hyperref[II.1.2.1-fr]{1.2.1}}
\egaTerme{Algèbre graduée sur un anneau gradué}{\hyperref[II.2.1.2-fr]{2.1.2}}
\egaTerme{Algèbre symétrique d'un $A$-module}{\hyperref[II.1.7.1-fr]{1.7.1}}
\egaTerme{$\mathcal{A}$-Algèbre symétrique d'un $\mathcal{A}$-Module}{\hyperref[II.1.7.4-fr]{1.7.4}}
\egaTerme{$\mathcal{O}_Y$-Algèbre graduée essentiellement réduite}{\hyperref[II.3.1.12-fr]{3.1.12}}
\egaTerme{$\mathcal{O}_Y$-Algèbre graduée intègre}{\hyperref[II.3.1.12-fr]{3.1.12}}
\egaTerme{Ample ($\mathcal{O}_X$-Module inversible)}{\hyperref[II.4.5.3-fr]{4.5.3}}
\egaTerme{Ample relativement à $f$, $f$-ample, $Y$-ample ($\mathcal{O}_X$-Module inversible)}{\hyperref[II.4.6.1-fr]{4.6.1}}
\egaTerme{Anneau gradué essentiellement intègre}{\hyperref[II.2.1.11-fr]{2.1.11}}
\egaTerme{Anneau gradué essentiellement réduit}{\hyperref[II.2.1.10-fr]{2.1.10}}
\egaTerme{Associé à un $\mathcal{B}$-Module (faisceau)}{\hyperref[II.1.4.3-fr]{1.4.3}}
\egaTerme{Associé à un $S$-module gradué (faisceau)}{\hyperref[II.2.5.3-fr]{2.5.3}}
\egaTerme{Associé à un $\mathcal{S}$-Module gradué (faisceau)}{\hyperref[II.3.2.2-fr]{3.2.2}}
\egaTerme{Associé à une $\mathcal{O}_S$-Algèbre ($S$-schéma)}{\hyperref[II.1.3.1-fr]{1.3.1}}

\egaTerme{Condition (TF), condition (TN)}{\hyperref[II.2.7.2-fr]{2.7.2}}
\egaTerme{Condition (TF), condition (TN)}{\hyperref[II.3.4.2-fr]{3.4.2}}
\egaTerme{Cône affine, cône projectif défini par une $\mathcal{O}_Y$-Algèbre graduée}{\hyperref[II.8.3.1-fr]{8.3.1}}
\egaTerme{Cône affine épointé, cône projectif épointé défini par une $\mathcal{O}_Y$-Algèbre graduée}{\hyperref[II.8.3.4-fr]{8.3.4}}
\egaTerme{Cône projetant affine, cône projetant projectif d'un préschéma $\Proj(\mathcal{S})$}{\hyperref[II.8.3.1-fr]{8.3.1}}
\egaTerme{Courbe algébrique sur un corps}{\hyperref[II.7.4.2-fr]{7.4.2}}
\egaTerme{Courbe algébrique complète}{\hyperref[II.7.4.7-fr]{7.4.7}}
\egaTerme{Critère de Serre}{\hyperref[II.5.2.1-fr]{5.2.1}}

\egaTerme{Déterminant d'un endomorphisme d'un module de type fini sur un anneau intègre}{\hyperref[II.6.4.2-fr]{6.4.2}}
\egaTerme{Déterminant d'un endomorphisme d'un module de type fini sur un anneau noethérien réduit}{\hyperref[II.6.4.7-fr]{6.4.7}}
\egaTerme{Déterminant d'un endomorphisme d'un $\mathcal{A}$-Module localement libre}{\hyperref[II.6.4.8-fr]{6.4.8}}
\egaTerme{Déterminant d'un endomorphisme d'un $\mathcal{O}_X$-Module de type fini sur un préschéma $X$ localement intègre}{\hyperref[II.6.4.9-fr]{6.4.9}}
\egaTerme{Déterminant d'un endomorphisme d'un $\mathcal{O}_X$-Module de type fini sur un préschéma $X$ localement noethérien et réduit}{\hyperref[II.6.4.10-fr]{6.4.10}}

\egaTerme{Éclaté (préschéma)}{\hyperref[II.8.1.3-fr]{8.1.3}}
\egaTerme{Entier (morphisme)}{\hyperref[II.6.1.1-fr]{6.1.1}}
\egaTerme{Entier (préschéma) au-dessus d'un préschéma}{\hyperref[II.6.1.1-fr]{6.1.1}}
\egaTerme{Entière ($\mathcal{O}_S$-Algèbre quasi-cohérente)}{\hyperref[II.6.1.2-fr]{6.1.2}}
\egaTerme{Entière (section d'une $\mathcal{A}$-Algèbre) sur $\mathcal{A}$}{\hyperref[II.6.3.1-fr]{6.3.1}}

\egaTerme{Faisceau fondamental d'un fibré projectif}{\hyperref[II.4.1.1-fr]{4.1.1}}
\egaTerme{Fermeture intégrale d'un faisceau d'anneaux $\mathcal{A}$ dans une $\mathcal{A}$-Algèbre}{\hyperref[II.6.3.2-fr]{6.3.2}}
\egaTerme{Fermeture intégrale d'un préschéma $X$ relativement à une $\mathcal{O}_X$-Algèbre}{\hyperref[II.6.3.4-fr]{6.3.4}}
\egaTerme{Fermeture projective d'un cône affine}{\hyperref[II.8.3.1-fr]{8.3.1}}

